\chapter{Venus writes Euclid in the sky}
\chaptertag{Three clocks, five alignments, and the Dresden Codex}
\label{app:venus}

Continued fractions are useful whenever two repeating motions almost return to
their starting arrangement at the same time.  Venus gives a particularly
visible example.  Small Fibonacci numbers describe a near-repeat of the modern
planetary motions, while related integers organize an exact calendar table in
the Dresden Codex.  We will build the astronomical vocabulary from scratch and
keep the approximate observation separate from the exact calendar arithmetic.

\section{Begin with two runners on circular tracks}

Imagine Earth and Venus as two runners circling the Sun.  Venus uses the inner
track and completes a lap sooner, so it repeatedly catches and passes Earth.
At the instant shown below, the Sun, Venus, and Earth lie on the same ray.

\begin{center}
\begin{tikzpicture}[scale=1.15]
  \draw[navy!28,line width=.9pt] (0,0) circle (1.25);
  \draw[navy!28,line width=.9pt] (0,0) circle (2.15);
  \draw[navy!22,dashed] (0,0)--(2.55,0);
  \fill[gold] (0,0) circle (4pt) node[below left=3pt,text=ink] {Sun};
  \fill[coral] (1.25,0) circle (4pt) node[below=4pt,text=ink] {Venus};
  \fill[teal] (2.15,0) circle (4.5pt) node[above=4pt,text=ink] {Earth};
  \draw[-{Stealth[length=2.2mm]},coral,line width=1.1pt]
    (120:1.25) arc[start angle=120,end angle=200,radius=1.25];
  \draw[-{Stealth[length=2.2mm]},teal,line width=1.1pt]
    (120:2.15) arc[start angle=120,end angle=180,radius=2.15];
  \node[coral,align=center] at (-1.55,-1.45) {shorter lap\\faster runner};
  \node[teal,align=center] at (-2.45,1.65) {longer lap\\slower runner};
  \node[text=muted,align=center] at (4.0,.05)
    {an alignment\\(a conjunction)};
  \draw[-{Stealth[length=2mm]},muted] (3.25,.05)--(2.35,.05);
\end{tikzpicture}
\end{center}

Here are the three time intervals we need.

\begin{center}
\small
\begin{tabularx}{\textwidth}{>{\bfseries}p{.16\textwidth}p{.30\textwidth}X}
\toprule
Symbol & Everyday meaning & Astronomical name and approximate value \\
\midrule
$E$ & time for Earth to make one Sun-centered lap
& Earth's \emph{sidereal period}: $365.256$ days \\
$V$ & time for Venus to make one Sun-centered lap
& Venus's \emph{sidereal period}: $224.7$ days \\
$S$ & time from one Earth--Venus alignment to the next
& Venus's \emph{synodic period}: $583.92$ days \\
\bottomrule
\end{tabularx}
\end{center}

``Sidereal'' means that a full lap is measured relative to the distant stars.
A ``conjunction'' is an alignment; in the alignment drawn above, Venus lies
between Earth and the Sun, so astronomers call it an \emph{inferior
conjunction}.  ``Synodic period'' simply means the waiting time until the same
kind of alignment occurs again.\footnote{For the modern mean values, see NASA,
\href{https://eclipse.gsfc.nasa.gov/OH/transit12.html}{\emph{2012 Transit of
Venus}} and \href{https://history.arc.nasa.gov/hist_pdfs/nasa_sp461.pdf}{NASA
SP--461, \emph{Pioneer Venus}}.  The real orbits are neither perfectly circular
nor perfectly uniform; rounded mean periods are all we need here.}

\section{The overtaking clock}

Earth completes $1/E$ laps per day, while Venus completes $1/V$.  Because Venus
is faster, it gains on Earth at the rate
\[
  \frac1V-\frac1E\quad\hbox{laps per day}.
\]
The next alignment occurs when Venus has gained exactly one whole lap.  Hence
\begin{equation}
  \frac1S=\frac1V-\frac1E.                         \label{eq:venus-beat}
\end{equation}
This is the same ``lapping time'' calculation used for two runners on a track.

Put $r=E/V$: it compares the length of an Earth lap in time with a Venus lap.
Multiplying \eqref{eq:venus-beat} by $E$ gives
\begin{equation}
  \frac SE=\frac1{r-1}=:T(r).                      \label{eq:venus-mobius}
\end{equation}
The measured ratio $E/V\approx1.6255$ has the simple, very accurate fraction
$13/8$ nearby.  If we replace the measured ratio temporarily by this ideal one,
then
\[
  T\!\left(\frac{13}{8}\right)=\frac85.
\]
In plain language:
\begin{itemize}
  \item during about eight Earth laps, Venus makes about thirteen laps; and
  \item in that time Venus gains $13-8=5$ laps on Earth, producing about five
  repeat alignments.
\end{itemize}
None of these counts is exact for the physical planets.  Their usefulness is
that $13/8$ is an excellent approximation with a tiny denominator.

\begin{checkpoint}[title={Why the golden ratio appears}]
The positive fixed point of the overtaking map $T(r)=1/(r-1)$ satisfies
\[
  r=\frac1{r-1},\qquad r^2-r-1=0,
\]
so it is $\varphi=(1+\sqrt5)/2$.  More generally,
\[
  T\!\left(\frac{F_{n+1}}{F_n}\right)
  =\frac{F_n}{F_{n-1}}.
\]
The calculation of one lapping time walks backward through consecutive
Fibonacci ratios.
\end{checkpoint}

\section{Why five alignments draw a star}

During one alignment-to-alignment interval, Earth travels about
$S/E\approx8/5$ laps.  One complete lap returns to the same direction, so only
the leftover $3/5$ of a lap changes the direction of the alignment ray.  The
five successive directions are therefore approximately
\[
  0,\quad \frac35,\quad \frac15,\quad \frac45,\quad \frac25
  \pmod 1.
\]
Connecting those directions in time order draws a five-pointed star.

\begin{center}
\begin{tikzpicture}[scale=1.2]
  \coordinate (v0) at (90:2);
  \coordinate (v1) at (162:2);
  \coordinate (v2) at (234:2);
  \coordinate (v3) at (306:2);
  \coordinate (v4) at (18:2);
  \draw[navy!30,line width=.8pt] (0,0) circle (2);
  \draw[navy!20,dashed] (0,0) -- (v0);
  \draw[coral,line width=1.8pt]
    (v0) -- (v3) -- (v1) -- (v4) -- (v2) -- cycle;
  \fill[gold] (0,0) circle (3pt) node[below=4pt,text=muted] {Sun};
  \foreach \v/\lab/\pos in {
    v0/$C_0$/above,
    v3/$C_1$/below right,
    v1/$C_2$/above left,
    v4/$C_3$/above right,
    v2/$C_4$/below left}
    \fill[teal] (\v) circle (2.7pt) node[\pos,text=ink] {\lab};
\end{tikzpicture}

\small\color{muted}Each $C_j$ marks the direction of an alignment, not the
path traveled by Venus.  The actual fifth return misses the starting direction
slightly.
\end{center}

The geometry of a perfectly regular pentagram contains the golden ratio
exactly.  The planetary star is only approximate: because $13/8$ is not the
exact period ratio, the alignment points drift and the whole pattern slowly
turns over repeated eight-year cycles.\footnote{The European Southern
Observatory's educational note
\href{https://www.eso.org/public/outreach/eduoff/vt-2004/Education/edu1app5.html}{\emph{The
Pentagram of Venus}} illustrates the five alignments and the small miss after
eight years.}

\section{The Euclidean algorithm inside \texorpdfstring{$13:8$}{13:8}}

The approximation contains the short Euclidean chain
\[
\begin{aligned}
  13&=1\cdot8+5, & 8&=1\cdot5+3,\\
   5&=1\cdot3+2, & 3&=1\cdot2+1,\\
   2&=2\cdot1.
\end{aligned}
\]
Consequently
\[
  \frac{13}{8}=[1;1,1,1,2]=[1;1,1,1,1,1],
  \qquad
  \frac85=[1;1,1,2]=[1;1,1,1,1].
\]
These are consecutive convergents of
$\varphi=[1;\overline1]$.  The orbit ratio $13/8$, the difference $5$, the
step $3/5$, and the pentagram are not independent coincidences: they are
different consequences of one good rational approximation.

\section{The Dresden Codex: an exact calendar model}

The modern calculation above starts with measured periods and obtains
approximations.  The Maya Venus Table in the Dresden Codex instead uses chosen
whole-day counts for calendrical bookkeeping.  In particular:
\begin{itemize}
  \item $584$ days is the table's rounded alignment-to-alignment Venus cycle;
  \item $365$ days is one \emph{Haab}, the Maya solar-year calendar; and
  \item $260$ days is one Tzolk'in, a repeating named-day cycle.
\end{itemize}

Five pages of the table present one $584$-day Venus cycle apiece.  The five
cycles make an eight-Haab block because
\begin{equation}
  5\cdot584=8\cdot365=2920.                         \label{eq:dresden-block}
\end{equation}
This equality is exact arithmetic inside the calendar model, not an assertion
that the physical Venus cycle is exactly $584$ days.\footnote{See the Saxon
State and University Library Dresden's description of the
\href{https://www.slub-dresden.de/en/explore/manuscripts/the-dresden-maya-codex/content}{Dresden
Maya Codex}, especially the Venus pages 24 and 46--50.  For the page structure
and cumulative day counts, see Akira Kaneko,
\href{https://classic.scielo.org.mx/scielo.php?pid=S0185-25742010000100002&script=sci_arttext}{``Identification
of Postclassic Maya Constellations from the Venus Pages of the Dresden
Codex,''} \emph{Estudios de Cultura Maya} \textbf{35} (2010).}

To make the $2920$-day block agree again with the $260$-day named-day cycle,
take a least common multiple.  Since $\gcd(2920,260)=20$,
\[
  \operatorname{lcm}(2920,260)=\frac{2920\cdot260}{20}=37960.
\]
At that common return,
\begin{equation}
  65\cdot584=104\cdot365=146\cdot260
  =13\cdot2920=37960.                              \label{eq:dresden-long}
\end{equation}

\begin{labbox}[title={The same numbers are doing different jobs}]
In the modern orbital approximation, $13$ counts Venus's Sun-centered laps in
about eight Earth laps.  In \eqref{eq:dresden-long}, $13$ counts complete
$2920$-day calendar blocks in a $37960$-day table.  The identity $13=8+5$ is a
delightful number-theoretic lens, but the manuscript alone does not show that
Maya astronomers formulated their table using Fibonacci numbers or continued
fractions.
\end{labbox}

Finally, rounding has a cost.  The modern mean alignment cycle is about
$583.92$ days, so the calendar value $584$ gains roughly
\[
  65(584-583.92)\approx5.2\text{ days}
\]
over $65$ cycles.  Scholars study correction schemes in the table.  Our exact
identity \eqref{eq:dresden-long} describes its integer framework, not a
perfectly uniform planet.

\begin{exercisebox}
\begin{enumerate}
  \item In your own words, distinguish a planet's Sun-centered lap time from
  the waiting time between two Earth--planet alignments.  Which are $E$, $V$,
  and $S$? \exlevel{Fluency}

  \item Derive \eqref{eq:venus-beat} by comparing lap rates.  Then derive
  \eqref{eq:venus-mobius} and verify that $T(13/8)=8/5$.
  \exlevel{Reasoning}

  \item Starting from a step of $3/5$ of a turn, list the first five positions
  modulo $1$.  Why does joining them chronologically make a star rather than
  the boundary of a pentagon? \exlevel{Reasoning}

  \item Run the Euclidean algorithm on $(13,8)$.  Read its quotients as the
  continued fraction of $13/8$, and explain why replacing the final $2$ by
  $1,1$ gives the same rational number. \exlevel{Reasoning}

  \item Factor $2920$ and $260$, compute their gcd and lcm, and recover every
  multiplier in \eqref{eq:dresden-long}. \exlevel{Fluency}

  \item Prove $T(F_{n+1}/F_n)=F_n/F_{n-1}$.  Show that the positive fixed point
  of $T$ is $\varphi$, and connect this with $[1;\overline1]$.
  \exlevel{Reasoning}

  \item Using $583.92$ days as the modern mean alignment cycle, compare the
  elapsed time for $5$, $65$, and $650$ cycles with the nominal $584$-day
  count.  Which error is negligible for a sketch but serious for a long
  calendar? \exlevel{Exploration}
\end{enumerate}
\end{exercisebox}
